\documentclass[conference]{IEEEtran}
\IEEEoverridecommandlockouts
\usepackage{cite}
\usepackage{amsmath,amssymb,amsfonts}
\usepackage{algorithmic}
\usepackage{graphicx}
\usepackage{textcomp}
\usepackage{xcolor}
\usepackage[brazil]{babel}

\def\BibTeX{{\rm B\kern-.05em{\sc i\kern-.025em b}\kern-.08em
    T\kern-.1667em\lower.7ex\hbox{E}\kern-.125emX}}  

\begin{document}

\title{Resenha do Artigo Hexagonal Architecture (Ports and Adapters)}

\author{
\IEEEauthorblockN{Artur Bomtempo Colen}
\IEEEauthorblockA{
Engenharia de Software \\
Pontifícia Universidade Católica de Minas Gerais (PUC Minas) \\
Belo Horizonte, Brasil \\
abcolen@sga.pucminas.br}
}

\maketitle

\begin{abstract}

Este trabalho apresenta uma análise do artigo sobre \textit{Hexagonal Architecture}, proposto por Alistair Cockburn. O autor introduz o padrão arquitetural conhecido como Ports and Adapters, cujo objetivo principal é separar a lógica de negócio das dependências externas, como interfaces de usuário, bancos de dados e serviços externos. O artigo discute problemas recorrentes em arquiteturas tradicionais e propõe uma abordagem mais flexível, que permite maior testabilidade, manutenção e evolução do sistema. Esta resenha tem como objetivo apresentar os principais conceitos do artigo, demonstrar o entendimento geral sobre a proposta e discutir sua aplicação prática no contexto de desenvolvimento de software moderno.

\end{abstract}

\begin{IEEEkeywords}
arquitetura de software, arquitetura hexagonal, ports and adapters, desacoplamento, design de software
\end{IEEEkeywords}

\section{Introdução}

A arquitetura de software desempenha um papel fundamental na construção de sistemas robustos e sustentáveis. No entanto, muitos sistemas são desenvolvidos com forte acoplamento entre a lógica de negócio e os componentes externos, como interfaces gráficas e bancos de dados, o que dificulta sua evolução ao longo do tempo.

O artigo analisado apresenta a arquitetura hexagonal como uma solução para esse problema. A proposta central é reorganizar a estrutura do sistema de forma que a lógica principal fique isolada e independente de tecnologias externas.

De acordo com o autor, um dos principais erros no desenvolvimento de software é permitir que decisões técnicas influenciem diretamente a lógica de negócio. Isso resulta em sistemas difíceis de testar, manter e modificar. A arquitetura hexagonal busca evitar esse problema ao introduzir uma separação clara entre o núcleo da aplicação e suas interações externas.

Além disso, o artigo propõe uma mudança de mentalidade: em vez de desenvolver sistemas orientados a tecnologias, deve-se desenvolver sistemas orientados ao domínio do problema.

\section{O Conceito de Ports and Adapters}

O conceito central da arquitetura hexagonal é baseado em portas (\textit{ports}) e adaptadores (\textit{adapters}). As portas representam interfaces que definem como o sistema se comunica com o ambiente externo, enquanto os adaptadores são responsáveis por implementar essas interfaces de acordo com uma tecnologia específica.

Conforme descrito no artigo (páginas 2 e 3) :contentReference[oaicite:0]{index=0}, a aplicação deve ser capaz de funcionar independentemente de como é acessada ou de onde os dados vêm. Isso significa que uma mesma funcionalidade pode ser utilizada por diferentes meios, como interface gráfica, API ou testes automatizados.

Os adaptadores atuam como intermediários, convertendo dados externos para um formato compreensível pela aplicação e vice-versa. Dessa forma, a lógica interna permanece intacta, enquanto os detalhes de implementação ficam isolados.

Essa abordagem promove um alto nível de desacoplamento, permitindo que componentes externos sejam modificados ou substituídos sem impactar o funcionamento do sistema.

\section{Estrutura da Arquitetura Hexagonal}

A arquitetura é representada visualmente como um hexágono para reforçar a ideia de múltiplas interfaces de interação, sem uma hierarquia rígida como nas arquiteturas em camadas.

De acordo com o diagrama apresentado no artigo (página 4) :contentReference[oaicite:1]{index=1}, a aplicação ocupa a posição central, enquanto os adaptadores ficam ao redor, conectados por meio das portas.

Diferentemente do modelo tradicional em camadas, onde há uma dependência sequencial entre níveis, a arquitetura hexagonal permite que diferentes adaptadores se conectem diretamente ao núcleo da aplicação, de forma independente.

Outro aspecto importante é a simetria entre entradas e saídas. O artigo enfatiza que tanto a entrada de dados quanto a saída devem ser tratadas de maneira uniforme, o que simplifica o modelo mental do sistema.

Essa estrutura também facilita a reutilização de componentes e a substituição de tecnologias, tornando o sistema mais flexível.

\section{Problemas que a Arquitetura Resolve}

O artigo destaca diversos problemas comuns em sistemas tradicionais que são mitigados pela arquitetura hexagonal.

Um dos principais problemas é o acoplamento entre lógica de negócio e infraestrutura. Quando regras de negócio são implementadas diretamente em interfaces ou acessos a banco de dados, o sistema se torna rígido e difícil de modificar.

Outro problema relevante é a dificuldade de realizar testes automatizados. Em arquiteturas tradicionais, muitas vezes é necessário utilizar banco de dados real ou interfaces completas para testar o sistema, o que torna o processo mais lento e complexo.

Conforme exemplificado no artigo (páginas 5 a 7) :contentReference[oaicite:2]{index=2}, a arquitetura hexagonal permite utilizar adaptadores simulados (\textit{mocks}), possibilitando testes isolados e mais eficientes.

Além disso, o artigo menciona a dificuldade de integração com diferentes sistemas como um desafio recorrente. A utilização de portas padronizadas facilita essa integração, pois define contratos claros de comunicação.

\section{Entendimento Geral do Artigo}

A partir da análise do artigo, é possível compreender que a arquitetura hexagonal representa uma evolução importante na forma de estruturar sistemas.

Um dos principais aprendizados é que a independência da lógica de negócio em relação às tecnologias externas é essencial para a longevidade do software. Sistemas que dependem fortemente de tecnologias específicas tendem a se tornar obsoletos mais rapidamente.

Outro ponto relevante é a importância dos contratos definidos pelas portas. Eles garantem que diferentes partes do sistema possam evoluir de forma independente, desde que respeitem as interfaces estabelecidas.

Também fica evidente que a arquitetura incentiva boas práticas de engenharia de software, como separação de responsabilidades, modularidade e testabilidade.

Por fim, o artigo reforça a ideia de que decisões arquiteturais devem ser tomadas com foco na manutenção e evolução do sistema, e não apenas na facilidade de implementação inicial.

\section{Aplicação Prática}

Na prática, a arquitetura hexagonal pode ser aplicada em diversos cenários do desenvolvimento de software.

Um exemplo comum é o desenvolvimento de sistemas web. A lógica de negócio pode ser implementada no núcleo da aplicação, enquanto diferentes adaptadores permitem acesso via REST, interfaces web ou testes automatizados.

Outro exemplo é a persistência de dados. Em vez de acoplar o sistema diretamente a um banco de dados específico, é possível definir uma porta para acesso a dados e criar adaptadores para diferentes tecnologias, como bancos relacionais ou NoSQL.

A arquitetura também é amplamente utilizada em sistemas que precisam de alta testabilidade. Com o uso de adaptadores simulados, é possível executar testes sem depender de infraestrutura externa, tornando o processo mais rápido e confiável.

Além disso, em ambientes corporativos, onde sistemas precisam se integrar com múltiplos serviços, a arquitetura hexagonal facilita essa integração ao estabelecer interfaces bem definidas.

Dessa forma, sua aplicação contribui para a construção de sistemas mais flexíveis, reutilizáveis e preparados para mudanças.

\section{Conclusão}

O artigo de Alistair Cockburn apresenta uma abordagem clara e eficaz para lidar com problemas comuns no desenvolvimento de software.

A arquitetura hexagonal se destaca por promover o desacoplamento entre a lógica de negócio e os elementos externos, permitindo maior flexibilidade, testabilidade e facilidade de manutenção.

O principal entendimento obtido é que sistemas bem projetados devem ser independentes de tecnologias específicas, focando em contratos e responsabilidades bem definidas.

Além disso, a arquitetura incentiva boas práticas e contribui para o desenvolvimento de software mais sustentável ao longo do tempo.

Portanto, a adoção desse padrão pode trazer benefícios significativos tanto em projetos acadêmicos quanto profissionais.

\begin{thebibliography}{00}

\bibitem{b1}
A. Cockburn,
``Hexagonal Architecture (Ports and Adapters),''
2005.

\end{thebibliography}

\end{document}
